\section{The published equivalence comparison for quadric
fibrations}\label{the-published-equivalence-comparison-for-quadric-fibrations}

\subsection{CC-XC19 --- Rational and numerical equivalence on the total
space}\label{cc-xc19-rational-and-numerical-equivalence-on-the-total-space}

Let \(f:X\to B\) be a flat dominant morphism of smooth projective
varieties over a finite field \(k=\mathbf F_q\), with every geometric
fibre a quadric. Suppose \(\dim B\le2\). All cycle groups and
equivalence relations in this statement have rational coefficients.

\textbf{Published theorem.} Vial (2013, Corollary 4.8, p.~1536) proves
that rational and numerical equivalence agree on \(X\). Equivalently,
for every codimension \(j\), the quotient map

\[
  \mathrm{CH}^{j}(X)_{\mathbf Q}
       \longrightarrow
  \mathrm N^{j}(X)_{\mathbf Q}
\tag{CC-XC.Q5}
\]

is an isomorphism. Here \(\mathrm N^j\) denotes cycles modulo numerical
equivalence. The hypotheses include singular quadric fibres; smoothness
is required of the total space and base, not of the fibration. There is
no restriction to odd characteristic in the stated theorem.

This is an application of a published result, not a new proof of that
result. The theorem and its complete printed proof have been inspected
in the publisher's version. Its proof uses the Chow decomposition of
Corollary 4.4 and the vanishing of numerically trivial rational Chow
groups in the relevant degrees on varieties of dimension at most two
over a finite field.

\textbf{Scope corollary.} For every finite extension \(k'/k\), every
prime \(\ell\ne\operatorname{char}k\), and every \(j\), rational,
\(\ell\)-adic homological and numerical equivalence agree on the
rational cycles of \(X_{k'}\). In particular, the standard conjecture
\(D\), that homological and numerical equivalence agree, holds on
\(X_{k'}\), and the cycle-class map on
\(\mathrm{CH}^{j}(X_{k'})_{\mathbf Q}\) is injective.

\textbf{Complete deduction of the scope corollary.} Flatness, dominance,
projectivity and smoothness are preserved by the finite field extension,
and the geometric fibres remain quadrics. Apply Vial's Corollary 4.8 to
\(f_{k'}\). Rational equivalence implies homological equivalence.
Homological equivalence implies numerical equivalence, because the
cycle-class map respects intersection products and degrees. These two
implications, together with equality of the first and third equivalence
relations, force equality of all three. Equivalently, the homologically
trivial subgroup lies in the numerically trivial subgroup of the
rational Chow group, which is zero by the published theorem. This proves
the asserted injectivity and \(D\) for every indicated \(\ell\).
\(\square\)

The conclusion concerns cycles on the specified total space and its
finite field extensions. Corollary 4.8 does not assert this comparison
on every power \(X^m\), or on every product with another variety. Those
products do not automatically satisfy its hypothesis of being a quadric
fibration over a base of dimension at most two. The result also does not
assert surjectivity of the cycle-class map onto Frobenius-invariant
cohomology, which is the separate Tate existence problem.

For cote 15, this closes the equivalence-comparison question on each
such \(X\), while retaining the exact motive-generation boundary in
CC-XC18. In particular, over a surface base the present comparison does
not supply abelian-sector membership of the base or of the residual
surface projector in the even-dimensional quadric case. Actual
membership of the total space still has precisely the conditions proved
in CC-XC18.

\subsection{Primary reference}\label{primary-reference}

Vial, C. (2013). Algebraic cycles and fibrations. \emph{Documenta
Mathematica, 18}, 1521--1553. https://doi.org/10.4171/DM/435 . Publisher
PDF: https://ems.press/content/serial-article-files/26216?nt=1 . Exact
locator: Corollary 4.8 and its complete proof, p.~1536;
rational-coefficient convention, p.~1522; preceding decomposition,
Corollary 4.4, p.~1535. Received 23 March 2012; revised 19 July 2013.
These locators use the published 1521--1553 pagination.
